\section{Программная реализация}
\label{sec:implementation}

\subsection{Расчётные модули}
\draft{Расчётная часть написана на C++: геометрия, сетки, свойства среды,
дискретизация и решатели разделены по модулям. GMRES и POUNDERS вызываются
через PETSc/TAO. Подготовка входа, экспорт XLSX и визуализация выполняются
в Python. T и Ro сохраняются на кровле, подошве и в среднем по слоям
на исходных координатах; для обратной задачи выводится также найденный поток.}

\subsection{Минимальная высота ячеек и перенос}
\draft{Нулевая мощность нового слоя заменяется высотой ячейки
$\varepsilon=0{,}05$ м; суммарная мощность $N$ сжатых ячеек равна
$N\varepsilon$. Число ячеек фиксировано на эпохе, раскрытие происходит
снизу вверх. На рис.~\ref{fig:growth} показаны четыре стадии по двум
первым конфигурациям Модели 1 при $x=1400$ м. Мощность исходного слоя сохраняется,
его подошва опускается вместе с накоплением нового материала.}
\reportfigure{04_growth.png}
  {Накопление слоя: (а)--(г) --- четыре последовательных момента.}
  {fig:growth}
\draft{При смене сетки T интерполируется по вертикали в двух соседних колонках,
затем по горизонтали. Для интегралов реакций перенос выполняется отдельно
по сохраняющимся слоям. Новый материал получает температуру поверхности
и нулевые интегралы.}

\subsection{Упрощение исходной геометрии}
\draft{Расчётная сетка отделена от геометрии входных данных.
Алгоритм Рамера--Дугласа--Пекера (RDP)~\cite{ramer1972,douglas1973}
удаляет промежуточные точки полилиний с заданным допуском;
контакты и области существования слоёв сохраняются.
На рис.~\ref{fig:rdp} показаны исходная и упрощённые границы Модели 2
на нескольких участках и одинаковый участок сетки.
Кривые с допусками 1 и 5 м смещены вниз для раздельного просмотра;
внутри каждого участка шаг смещения одинаков для обеих кривых.}
\reportfigure{05_rdp.png}
  {Модель 2: (а)--(в) --- участки границы; (г), (д) --- сетка
  без упрощения и с допуском 5 м.}
  {fig:rdp}
\draft{При допуске 5 м число ячеек сокращается с 120335 до 3943.
Отличия средних по слоям T и Ro от подробного расчёта составляют
1,31 и 1,12\%, локальное отличие Ro --- до 6,91\%
(табл.~\ref{tab:rdp}). Эталон для Модели 2 не опубликован;
сравнение выполнено с расчётом без упрощения. Геометрия и сетка меняются
одновременно; время измерено однократно.}
\begin{table}[htbp]
\centering\small
\caption{Модель~2: упрощение геометрии и относительные отличия решения.}\label{tab:rdp}
\begin{tabular}{lrrrrr}
\toprule
Допуск, м & Ячеек & Время, с & $E_T$, \% & $E_{\Ro}$, \% & Макс. Ro, \% \\
\midrule
0 & 120335 & 205,41 & 0,00 & 0,00 & 0,00 \\
1 & 5288 & 7,79 & 0,56 & 0,92 & 3,48 \\
5 & 3943 & 4,95 & 1,31 & 1,12 & 6,91 \\
10 & 3723 & 3,62 & 1,32 & 1,18 & 6,90 \\
25 & 2955 & 2,72 & 1,36 & 1,11 & 7,49 \\
\bottomrule
\end{tabular}
\end{table}


\subsection{Хранение зрелости и табулирование}
\draft{T хранится в центрах ячеек, интегралы реакций --- в $N+1$ точке
на горизонтальных границах колонки. На внутренней границе температура
определяется с учётом тепловых сопротивлений:}
\begin{equation}
 T_\sigma=\frac{r_NT_P+r_PT_N}{r_P+r_N},\qquad
 r_P=\frac{d_{P\sigma}}{\lambda_P},\quad
 r_N=\frac{d_{N\sigma}}{\lambda_N}.
\end{equation}
\draft{На поверхности используется заданная температура, на подошве ---
температура нижней ячейки и заданный поток. Моменты $t_{\mathrm{dep}}$
подготовлены для каждой точки расчёта интегралов.
При первом обращении строится таблица 8192 значений на $\chi\in[20,200]$;
между ними применяется линейная интерполяция:}
\begin{equation}
 H(\chi)=\chi^2e^\chi G(\chi),\qquad
 \Phi_i(T)=\frac{T^2}{\beta_i}e^{-\beta_i/T}H(\beta_i/T).
\end{equation}
\draft{Проверка выполнена на одномерном охлаждении слоя мощностью 4000 м
с аналитической температурой}
\begin{equation}
 T_*(z,t)=283{,}15+200\exp\!\left(-\frac{\lambda k^2}{C_0}t\right)\sin(kz),
 \qquad k=\frac{\pi}{8000\,\text{м}}.
\end{equation}
\draft{Заданы $\lambda=2$ Вт/(м$\cdot$К), $C_0=3\cdot10^6$ Дж/(м$^3\cdot$К).
Длительность --- 1 млн лет, 200 шагов, 81 точка по глубине.
Эталон Ro получен независимой квадратурой по этой температуре;
удвоение числа интервалов меняет Ro менее чем на $3{,}2\cdot10^{-7}\%$.
Максимальная ошибка расчёта с линейной T внутри шага --- 0,00591\%;
вклад таблицы относительно интеграла по той же линейной T ---
$1{,}4\cdot10^{-8}\%$.}
\reportfigure{06_easy_ro.png}
  {(а) --- время вычисления интегралов; (б) --- ошибка Ro в одномерной задаче.}
  {fig:easy-speed}
\draft{На панели (а) использованы 65536 точек и 20 реакций,
пять повторов; показаны медианы и диапазоны времени.
Прямая первообразная вычисляется через интегральную экспоненту.
Таблица ускоряет вычисление интегралов в 6,55 раза;
это сравнение ядра вычисления интегралов, а не всего прямого расчёта.}

\subsection{Повторные расчёты и линейное решение}
\draft{При оптимизации геометрия и описания сеток эпох переиспользуются;
стационарное начальное поле пересчитывается для каждого потока.
Для MPFA структура матрицы сохраняется внутри нелинейного решения,
обновляются её значения. Векторы PETSc переиспользуются.
Промежуточные поля при подборе не записываются.}
\draft{Сравнение GMRES(100)/ILU(2) и UMFPACK~\cite{suitesparse} приведено
в табл.~\ref{tab:solvers}; распределение времени --- в табл.~\ref{tab:stages}.
В первой таблице проценты отсчитываются от времени UMFPACK
для соответствующего этапа; во второй --- от полного времени расчёта каждой модели.
Условия: Intel Core i7-12700H, один поток, GCC 15.2, оптимизированная сборка,
PETSc 3.24.4; $\Delta x=200$ м, $\Delta z=50$ м, $\Delta t=1$ млн лет,
$\varepsilon=0{,}05$ м, допуск RDP 5 м. Относительные допуски линейного
и нелинейного решений --- $10^{-8}$. Указаны медианы трёх запусков;
Python-подготовка и экспорт XLSX не включены.
Различия температур между решателями не превышают $3\cdot10^{-11}\,{}^\circ$C.}
\begin{table}[htbp]
\centering\small
\renewcommand{\arraystretch}{1.12}
\caption{Время расчёта в секундах и процентах от времени UMFPACK для того же этапа и модели.}\label{tab:solvers}
\begin{tabular}{ll D{,}{,}{2.3} D{,}{,}{3.1} D{,}{,}{2.3} D{,}{,}{3.1}}
\toprule
\multirow{2}{*}{Модель} & \multirow{2}{*}{Решатель} & \multicolumn{2}{c}{Линейное решение} & \multicolumn{2}{c}{Прямой расчёт} \\
\cmidrule(lr){3-4}\cmidrule(lr){5-6}
\multicolumn{1}{c}{} & \multicolumn{1}{c}{} & \multicolumn{1}{c}{с} & \multicolumn{1}{c}{\%} & \multicolumn{1}{c}{с} & \multicolumn{1}{c}{\%} \\
\midrule
\multirow{2}{*}{Модель~1} & GMRES/ILU(2) & 1,160 & 24,7 & 2,291 & 39,2 \\
 & UMFPACK & 4,698 & 100,0 & 5,839 & 100,0 \\
\addlinespace[3pt]
\multirow{2}{*}{Модель~2} & GMRES/ILU(2) & 2,471 & 27,3 & 4,995 & 43,1 \\
 & UMFPACK & 9,062 & 100,0 & 11,586 & 100,0 \\
\addlinespace[3pt]
\multirow{2}{*}{Модель~3} & GMRES/ILU(2) & 16,049 & 69,2 & 22,486 & 76,2 \\
 & UMFPACK & 23,178 & 100,0 & 29,511 & 100,0 \\
\bottomrule
\end{tabular}
\end{table}

\begin{table}[htbp]
\centering\small
\renewcommand{\arraystretch}{1.12}
\caption{Время этапов и доля полного прямого расчёта.}\label{tab:stages}
\begin{tabular}{l D{,}{,}{2.3} D{,}{,}{3.1} D{,}{,}{2.3} D{,}{,}{3.1} D{,}{,}{2.3} D{,}{,}{3.1} }
\toprule
\multirow{2}{*}{Этап} & \multicolumn{2}{c}{Модель~1} & \multicolumn{2}{c}{Модель~2} & \multicolumn{2}{c}{Модель~3} \\
\cmidrule(lr){2-3}\cmidrule(lr){4-5}\cmidrule(lr){6-7}
\multicolumn{1}{c}{} & \multicolumn{1}{c}{с} & \multicolumn{1}{c}{\%} & \multicolumn{1}{c}{с} & \multicolumn{1}{c}{\%} & \multicolumn{1}{c}{с} & \multicolumn{1}{c}{\%} \\
\midrule
Стационарный старт & 0,002 & 0,1 & 0,002 & 0,0 & 0,010 & 0,0 \\
Геометрия и сетка & 0,027 & 1,2 & 0,050 & 1,0 & 0,170 & 0,8 \\
Перенос & 0,009 & 0,4 & 0,014 & 0,3 & 0,043 & 0,2 \\
Сборка & 0,929 & 40,5 & 1,751 & 35,1 & 5,235 & 23,3 \\
Линейное решение & 1,160 & 50,6 & 2,471 & 49,5 & 16,049 & 71,4 \\
EASY\%Ro & 0,093 & 4,1 & 0,165 & 3,3 & 0,528 & 2,3 \\
Вывод & 0,068 & 3,0 & 0,541 & 10,8 & 0,423 & 1,9 \\
Прочее & 0,004 & 0,2 & 0,001 & 0,0 & 0,028 & 0,1 \\
\midrule
Всего & 2,291 & 100,0 & 4,995 & 100,0 & 22,486 & 100,0 \\
\bottomrule
\end{tabular}
\end{table}

